\chapter{Complete workflows}
\label{ch:workflows}

This chapter assembles the components of the previous chapters into the two workflows of an experiment: the
reconstruction of the data and the simulation used to interpret them. The worked example is the a1975 experiment, a
$^{16}$C beam at 184~MeV on a 300~torr D$_2$ target in a 2.85~T field, whose macros live in
\code{macro/Unpack\_HDF5/a1975/} (data) and \code{macro/Simulation/ATTPC/16C\_dt/} (simulation of the
$^{16}$C(d,t)$^{15}$C channel). The analysis downstream of the kinematics caches is plain ROOT and is kept in a
separate repository; this chapter stops at the reconstructed and fitted events.

\section{Data: from HDF5 files to fitted tracks}

The data reconstruction runs in two passes, so that the slow fit can be repeated with different settings on the same
pattern-recognition output.

\subsection{Pass 1: unpacking, PSA, cleaning, pattern recognition}

\code{D2\_UKF/unpackReco\_multifit.C} builds one \code{FairRunAna} with the parameter file, the pad map and four
tasks. The production script \code{D2\_UKF/dv\_dt.sh} calls it per run with the multi-fit PSA, a threshold of 40 and
the HDBSCAN track finder (the macro's default finder is \code{AtTrackFinderTC}):
\begin{verbatim}
root -l -b -q "unpackReco_multifit.C(\"$r\",0,false,\"$REC\",\"$H5\",false,true,
               \"multifit\",0,40,\"hdbscan\",0,3,0,0.1,\"$PAR\")"
\end{verbatim}
The tasks, as set up in the macro:
\begin{verbatim}
auto fAtMapPtr = std::make_shared<AtTpcMap>();
fAtMapPtr->ParseXMLMap(mapDir.Data());              // scripts/ANL2023.xml
fAtMapPtr->GeneratePadPlane();

auto unpacker = std::make_unique<AtHDFUnpacker>(fAtMapPtr);
unpacker->SetInputFileName(inputFile.Data());
unpacker->SetNumberTimestamps(2);
unpacker->SetBaseLineSubtraction(true);
auto unpackTask = new AtUnpackTask(std::move(unpacker));

auto p = new AtPSAMultiFit();                       // multi-peak pulse fit
p->SetThreshold(thr);
p->SetPeakingTime(0.720);
p->SetMaxPeaks(4);
p->SetMinSeparation(4);
AtPSAtask *psaTask = new AtPSAtask(p);
psaTask->SetOutputBranch("AtEventH");

auto cleaner = std::make_unique<AtTools::DataCleaning::AtDirDeDxCleaner>();
cleanTask = new AtDataCleaningTask(std::move(cleaner));    // -> AtEventClean

auto h = std::make_unique<AtPATTERN::AtTrackFinderHDBSCAN>();
h->SetMinClusterSize(hdMcs);                        // 0 = adaptive
h->SetMinSamples(hdMs);                             // 3
h->SetClusterSelectionEpsilon(10.0);
h->SetJoinMethod("mover");
h->SetMinClusterSizeJoin(15);
h->SetCircleOverlapRatio(0.25);
h->SetMotionGapTol(40);
h->SetMotionAngleTol(35);
AtPRAtask *praTask = new AtPRAtask(std::move(h));   // -> AtPatternEvent
\end{verbatim}
The cleaning step removes scattered background before track finding; the comment in the macro records its effect on
run 0016 (clean-proton yield from 5 to 20\,\%, median $\chi^2/\mathrm{ndf}$ of the fits from 206 to 59). The output,
\code{run\_XXXX\_multifit\_reco.root}, holds the hit cloud and the track candidates. The script writes a completion
marker only when the macro reports success, the output exceeds a size floor and the log has no segmentation fault,
and a failed stage deletes its own partial output. Chapter~\ref{ch:unpacking} describes the unpacker and the PSA,
Chapter~\ref{ch:patternrecognition} the track finders.

\subsection{Pass 2: PID and fit}

\code{D2\_UKF/fitGenfitter\_a1975\_deuterium.C} reads the pass-1 file and adds two tasks:
\begin{verbatim}
AtFitterTask *fitterTask = new AtFitterTask(std::move(fitter));  // AtGenfitter
fitterTask->SetInputBranch("AtPatternEvent");
fitterTask->SetOutputBranch("AtTrackingEvent");
AtPIDTask *pidTask = new AtPIDTask();
pidTask->SetInputBranch("AtPatternEvent");
pidTask->SetOutputBranch("AtPIDEvent");
run->AddTask(pidTask);
run->AddTask(fitterTask);
\end{verbatim}
The fitter takes the particle hypothesis, the material model and an optional PID gate (a JSON file of graphical cuts)
from the macro arguments. The production script \code{D2\_UKF/prod\_dt\_catima.sh} runs it for the triton hypothesis
on every run in parallel, checks each log for the banners that prove the intended energy-loss model was active, and
writes a \code{.done} marker per run so that a rerun skips what is complete. Chapter~\ref{ch:fitting} explains every
option; Chapter~\ref{ch:pid} the PID configuration.

\subsection{Reading the fitted events}

The fit output pairs each fitted track with its PID result through the track ID. The pattern used in
\code{D2\_UKF/ex\_dt\_a1975.C}:
\begin{verbatim}
tu->SetBranchAddress("AtTrackingEvent", &te);
tu->SetBranchAddress("AtPIDEvent", &pe);
for (Long64_t i = 0; i < tu->GetEntries(); ++i) {
   tu->GetEntry(i);
   if (pe->GetEntries() == 0 || te->GetEntries() == 0) continue;
   auto *pidev = (AtPIDEvent *)pe->At(0);
   auto *ev = (AtTrackingEvent *)te->At(0);
   std::map<int, AtFittedTrack *> fmap;
   for (auto &ft : ev->GetFittedTracks())
      if (ft) fmap[ft->GetTrackID()] = ft.get();
   for (auto &sr : pidev->GetSpyral()) {
      if (!sr.valid) continue;
      auto it = fmap.find(sr.trackID);
      if (it == fmap.end()) continue;
      auto &k = it->second->GetKinematicsXtr();   // at the vertex
      // ... k.kineticEnergy, k.theta, k.phi
   }
}
\end{verbatim}
Looping over the PID results rather than over the fitted tracks matters: the fit trees can repeat an
\code{AtTrackingEvent} in the next entry, and the copy carries an empty \code{AtPIDEvent}, so this loop skips it while a
loop over \code{GetFittedTracks()} would count the event twice.

\section{Simulation: from the reaction to reconstructed tracks}

\subsection{Generation}

\code{C16\_dt\_sim.C} sets up \code{FairRunSim} with Geant4, the AT-TPC detector filled with the D$_2$ medium, a
uniform field and two generators: \code{AtTPCIonGenerator} for the beam and \code{AtTPC2Body} for the reaction (beam,
target, heavy residual and light ejectile, in that order). The script \code{run\_states\_dt.sh} runs it once per
final state of $^{15}$C and records a marker only when the log shows both the requested seed and successful
completion:
\begin{verbatim}
root -b -q -l "C16_dt_sim.C($NEV,2.,70.,\"TGeant4\",-28.5,
                            \"$OUT/${J}_sim.root\",$EX,$SEED)"
grep -q "RNG seed requested: $SEED" "$OUT/${J}_gen.log" || echo NO_SEED
grep -q "Macro finished successfully" "$OUT/${J}_gen.log" || echo NOT_DONE
\end{verbatim}
The field is given in the data convention ($-28.5$~kG) so that simulated tracks turn like the data
(Chapter~\ref{ch:conventions}).

\subsection{Digitisation and reconstruction}

\code{run\_reco\_C16dt.C} reads the Monte Carlo points and runs the digitisation (\code{AtClusterizeTask},
\code{AtPulseTask}) followed by the same reconstruction chain as the data production: multi-fit PSA, the
\code{AtDirDeDxCleaner} cleaning and the HDBSCAN finder. Its parameter file is a copy of the data's
(\code{ATTPC.a1975\_deuterium\_dv1104\_g35000.par}, produced by \code{make\_gain\_pars.sh}) with only the gain changed,
so that the drift velocity, \code{TBEntrance} and \code{ZPadPlane}, hence the $z$ scale, are identical. Two choices in it carry the lesson of this
example:
\begin{itemize}
\item \emph{Gain.} The gain in the parameter file sets how many electrons reach a pad, hence how many hits each track
  keeps above the PSA threshold. It was matched to the data on the number of hits per track (gain 35\,000, ratio 0.99
  to run 0016); the header of \code{run\_reco\_batch\_dt.sh} documents the scan.
\item \emph{Beam hole.} The real detector does not read out the region around the beam; the simulation inhibits the
  same pads (\code{AtBeamHole.h}), otherwise the clustering merges the beam into the ejectile tracks.
\end{itemize}
The simulation and the data then go through \code{fitGenfitter\_a1975\_deuterium.C} with the same settings, which is
what makes an acceptance measured on the simulation applicable to the data.

\section{Running at scale}

\begin{itemize}
\item Split productions by run (data) or by seed (simulation) and run several jobs in parallel; each simulated job
  needs a distinct seed.
\item Write logs per job and check them for positive evidence of the intended configuration (the seed, the
  parameter file, the energy-loss banner), not only for the absence of errors.
\item Write a completion marker only after those checks; a file that exists is not a job that finished.
\item Start long jobs with \code{nohup} from a script, and never edit a script while it is running: bash reads it
  incrementally.
\item Record the exact command and inputs of every product, so that it can be regenerated and so that a cached result
  can be checked against its inputs.
\end{itemize}
